\section{Hạn chế và hướng phát triển}

\subsection{Hạn chế hiện tại}
Kế thừa các giới hạn đã nêu ở Chương 1 (LI-1, LI-2, LI-3), sau quá trình
hiện thực trong 6 tuần, nhóm ghi nhận thêm các hạn chế cụ thể:
\begin{itemize}
    \item \textbf{Phụ thuộc phần cứng:} Hệ thống phụ thuộc vào thiết bị
    POS vật lý và đầu đọc QR/RFID/NFC thật để kiểm thử đầy đủ; môi
    trường thử nghiệm hiện tại chủ yếu mô phỏng bằng phần mềm.
    \item \textbf{Độ trễ đồng bộ mạng:} Đúng như giới hạn LI-2 đã nêu,
    dashboard vận hành có thể không phản ánh dữ liệu 100\% real-time cho
    tới khi giao dịch ngoại tuyến được đồng bộ hết.
    \item \textbf{Quy mô kiểm thử:} Do giới hạn hạ tầng miễn phí (Render
    free tier, Vercel hobby tier), nhóm chưa kiểm thử được ở quy mô
    5.000 hành khách đồng thời như mục tiêu phi chức năng đã đặt ra.
    \item \textbf{Module BLE:} Chưa được hiện thực trong phiên bản này.
\end{itemize}

\subsection{Hướng phát triển}
\begin{itemize}
    \item Hiện thực module định vị hành khách qua thiết bị đeo BLE (đã
    thiết kế sẵn ở đặc tả yêu cầu Chương 3, mục 3.11) cho phiên bản kế
    tiếp.
    \item Tích hợp cổng thanh toán trực tuyến (VNPAY/payOS/Stripe) để hỗ
    trợ hành khách nạp tiền trước vào tài khoản trên tàu qua ứng dụng
    mobile.
    \item Bổ sung phân tích dữ liệu lớn (dashboard nâng cao) để dự đoán
    nhu cầu hoạt động/tham quan bờ dựa trên lịch sử các chuyến đi trước.
    \item Kiểm thử tải (load testing) ở quy mô lớn hơn bằng công cụ
    JMeter/Locust trên hạ tầng cloud trả phí để xác nhận đầy đủ các chỉ
    tiêu phi chức năng đã đặt ra.
    \item Phát hành chính thức ứng dụng Mobile lên App Store và Google
    Play sau khi hoàn tất kiểm thử bảo mật và tuân thủ chính sách của
    hai nền tảng.
\end{itemize}